\section{Permuting binary addresses on tapes}\label{sec:permutations}

Consider a stream of records with two binary address fields $X,Y$.
One step of the Walsh algorithm sends the record at $(X,Y)$ to
$(X,Y\oplus X)$, where $\oplus$ denotes bitwise addition modulo two.
For three-bit fields, for instance, the record at $(101,011)$ goes to
$(101,110)$. Rotating $Y$ by five modulo eight
gives $000$ in the target field because of the carry. Separately flipping every
bit of $Y$ would require a scan for each bit. We group bit flips into
a small number of shifts and repair the records on which a carry
crosses between groups. This gives a cost governed by the width of
the two fields, which decreases at each level of the Walsh recursion.

Later, the conversion to Chinese remainder coordinates permutes small
blocks of address bits according to other bits of the address. In that
application the available budget depends on the width of the whole
address. We derive both bounds from Theorem~\ref{thm:interchange} and
the bit packing of \cite[Section~5.2]{Tape}.

Fix $1/2<\rho<1$. A stream of $M$ records, each of $R$ bits, has bit
volume $V=MR$. Let $k\ge2$ bound the address width up to a constant
factor, and assume $R=\Theta(k)$ with fixed constants. Scanning the
address counters once per record thus costs $O(V)$. All shapes have
complete ranges in the binary fields on which we operate. A leading
row range can have arbitrary positive integer length, and the spectator
ranges between active fields may be combined. An address permutation
moves whole records and requires no information about their contents.

The elementary stream operations of \cite[Lemma~3]{Tape} move a
single address bit to any position, or permute the four values of any
two specified bits, in $O(V)$ steps. The bounds include the binary
counters, rewinds and cleanup. Descriptors of the fields are supplied
in binary. Throughout this section, the tape count and alphabet are
fixed, the constants in the block widths and polynomial preparation
costs are fixed, and the time bounds include this preparation in the
top parameter $k$ even when the active fields are short. The space
bound in each case is $O(V)+k^{O(1)}$ visited cells. Interchanges
take linear space, and a fixed number of streams suffices for the
scans and sorts below. Adding an $O(k)$-bit address to a record
increases its size by a constant factor.

\subsection{Parallel flips in two address fields}

\begin{lemma}[Bulk XOR between address chunks]\label{lem:packed-xor}
Let $X,Y,Z$ be disjoint complete binary chunks of a common width
$L=O(k)$, occurring in any order with arbitrary intervening spectator
fields. For each target $j$ in a specified set of positions in $Y$,
choose a source position $i(j)$ in $X$. The simultaneous changes
\[
                 Y_j\longleftarrow Y_j\oplus X_{i(j)}
\]
take $O\bigl(V(1+L^\rho)(1+\log k)\bigr)+k^{O(1)}$ fixed multitape
steps. Every other address bit, including all of $Z$, is restored.
\end{lemma}

\begin{proof}
Choose a guard length $G-1$, with $G=C\lceil\log_2(2k)\rceil$ and a
fixed integer $C$ to be chosen below. For one pass, select disjoint
segments of $G$ consecutive bits in $Y$ and the corresponding segments
in $Z$. Write $j$ for the low position of a segment, numbering from bit
zero. We will flip $Y_j$ according to a Boolean control $c_j$ and
restore all other bits of both segments. For this pass we may use
any unchanged controls preceding both chunks, provided the functions
$c_j$ can be computed in $k^{O(1)}$ steps from descriptions of size
$k^{O(1)}$. For $c_j=X_{i(j)}$, interchange equal-width chunks to put
$X$ first and undo the change of layout at the end.

\emph{A shift acting on several segments.}
For two unrestricted integer coordinates $u,w$, let $[u]_h$ indicate
$u\equiv h\bmod2$. The following eight shifts use the current coordinate
values in each assignment, in row order from left to right
\cite[Lemma~13, Equations~(20)--(21)]{Tape}.
\[
\begin{array}{ll}
 u\leftarrow u-c[w]_0 & w\leftarrow w-c[u]_0\\
 u\leftarrow u+c[w]_0 & w\leftarrow w+c[u]_0\\
 w\leftarrow w+c[u]_1 & u\leftarrow u+c[w]_1\\
 w\leftarrow w-c[u]_1 & u\leftarrow u-c[w]_1 .
\end{array}
\]
For $c=0$ these shifts fix both integers. For $c=1$, on each
translated square $\{2a,2a+1\}\times\{2b,2b+1\}$ the first four
assignments cycle the low-bit pairs $(00\ 01\ 10)$ and the last four
cycle $(01\ 10\ 11)$. Their composition flips the low bit of $u$
and restores $w$.

Apply the shifts at all selected segments of $Y,Z$ at once. For
example, one assignment with target $Y$ is the single rotation
\[
      Y\longleftarrow
          Y\mathbin\pm\sum_j2^j c_j[Z_j=h]\pmod{2^L}.
\]
Interchange $Y$ and $Z$ if necessary to place the target last among
these chunks. For every prefix preceding the target, compute the shift
amount from the control bits and the parity bits of the other chunk.
The $2^L$ target values occur consecutively as blocks with a common
suffix length. Split the run at the wrap point and concatenate its two
pieces in rotated order \cite[Lemma~4]{Tape}, keeping each suffix
block intact. Undo the interchange. Each prefix accounts for at least
$2^L\ge2^G$ records whenever there is a segment. Choosing $C$
sufficiently large covers the polynomial cost of calculating one offset
per prefix. With a fixed number of field counters and
Theorem~\ref{thm:interchange}, all eight rotations cost
$O(V(1+L^\rho))+k^{O(1)}$ steps.

\emph{Repairing carries.}
The upper $G-1$ bits of each selected segment form its guard.
Declare an address exceptional if any guard in $Y$ or $Z$ has one of
its first ten or last ten values. At every other address the integer
in each segment stays inside its $G$ bits throughout the eight shifts,
since its displacement is at most eight. Thus the performed permutation
$S$ agrees there with the desired controlled flips $D$. The map $D$
preserves every guard. Both $S$ and $D$ are bijections and agree on the
complement of the exceptional set $\mathcal B$, so they both preserve
$\mathcal B$. The remaining permutation $DS^{-1}$ need only be applied
to records now in exceptional slots.

With $s$ segment pairs the exceptional fraction is at most
$80s2^{-G}$, also within any fixed row, since each guard assumes all
$2^{G-1}$ values equally often. A scan identifies these records in
$O(k)$ steps per address using prepared masks for the guards. It
suffices to read their four low bits and test whether the higher bits
are all zero or all one. Extract each exceptional record and attach
its destination rank. Compute this rank from its current address by
reversing the eight rotations and then performing the prescribed bit
flips. This takes $k^{O(1)}$ steps per extracted record, including
address arithmetic in the rectangular layout. Sort these tagged
records by destination with $O(k)$ stable binary radix passes, then
merge them into the exceptional slots in stream order. Beyond setup
and full scans, this costs
\[
             |\mathcal B|\bigl(k^{O(1)}+O(kR)\bigr).
\]
There are $s=O(k)$ segment pairs. After fixing the polynomial bounds
for the controls, offsets and rank calculations, take $C$ large enough
that this cost is $O(V)$. This proves an exact pass even when all
auxiliary address values are occupied.

\emph{Covering the targets.}
Partition the target positions $j$ with $j+G\le L$ by their residues
modulo $G$. Within each class the segments are disjoint, so one pass
with $c_j=X_{i(j)}$ suffices. At most $G-1$ high positions have too
little room for a guard. Treat each by the elementary two-bit operation
on $X_{i(j)},Y_j$. This also handles the entire case $L<G$. There are
$O(1+\log k)$ passes in all.
\end{proof}

\subsection{Permuting positions and blocks}

\begin{corollary}[Permutations of address positions]\label{cor:bit-positions}
An explicitly specified permutation of the $k$ binary address positions
of $2^k$ records takes
$O_\rho(V(1+k)^\rho(1+\log k))$ deterministic fixed multitape steps.
\end{corollary}

\begin{proof}
A permutation is a product of two involutions. On each cycle, number
positions modulo its length and use the reflections $i\mapsto-i$ and
$i\mapsto1-i$ in that order. It suffices to perform the disjoint pairs
of exchanges in one involution. Split the address into three chunks
of width $\lfloor k/3\rfloor$ and at most two remaining bits. Handle
the pairs touching the remaining bits individually. Between any two
chunks perform all exchanges with three simultaneous XORs, using
the third chunk as the auxiliary chunk in Lemma~\ref{lem:packed-xor}.
Within a single chunk, assign a distinct temporary position $t$ in
another chunk to each exchange $(i,j)$. The three cross-chunk exchanges
$(i,t),(j,t),(i,t)$ effect $(i,j)$ and restore $t$. This uses a
bounded number of applications of the lemma. The $k^{O(1)}$ steps to
prepare the cycle decompositions and position lists are absorbed by
the full $2^k$-record volume.
\end{proof}

The Chinese remainder layout in Section~\ref{sec:integers} needs to
rotate the values of each of several short fields by amounts determined
by other address bits. The next lemma allows an arbitrary permutation
table for each such field, including a different table for every
assignment of the controls.

\begin{lemma}[Controlled permutations on small blocks]\label{lem:controlled-blocks}
On a full $k$-bit address space, take pairwise disjoint target blocks of
$O(\log(2k))$ specified bit positions each. Fix $\delta>0$ and a common
set of at most $(1-\delta)k$ control positions disjoint from the targets.
At each block prescribe a permutation of its values depending on the
control bits. At any control setting its table must be generable in
$k^{O(1)}$ tape steps, with uniform polynomial bounds on this computation
and its description. Applying all block permutations, or all their
inverses, to the records takes
\[
 O\bigl(V(1+k)^\rho(1+\log k)^3\bigr)+k^{O(1)}
\]
deterministic fixed multitape steps.
\end{lemma}

\begin{proof}
We first perform controlled flips on a batch of $b$ individual targets,
where $b\le c_\delta k/G$ for a sufficiently small $c_\delta>0$ and
$G$ is the guard parameter above. By Corollary~\ref{cor:bit-positions},
place all controls first and the targets at the low positions of $b$
disjoint $G$-bit segments of a chunk $Y$. Reserve a further $bG$
positions for $Z$. The two chunks fit in the at least $\delta k$ free
positions. Use the guarded pass from Lemma~\ref{lem:packed-xor} and
undo the rearrangement. Including the two position permutations, the
cost of one batch is $O(V(1+k)^\rho(1+\log k))+k^{O(1)}$. Thus a
layer of flips on all targets costs
$O(V(1+k)^\rho(1+\log k)^2)+k^{O(1)}$. Bounded $k$ can always be
handled by a direct split and merge.

For arbitrary block permutations we also need the other bits of each
active block as controls. Batch the blocks so that each batch uses at most
$\delta k/2$ positions and at most $c_{\delta/2}k/G$ blocks.
Greedily filling batches yields $O(1+\log k)$ batches. For large $k$,
a width-limited batch occupies $\Omega_\delta(k)$ positions since the
next block has only $O(\log k)$ bits. A count-limited batch contains
$\Omega_\delta(k/G)$ blocks.

Ignore empty blocks. On a $d$-bit block with $d\ge1$, any permutation of the $2^d$ values is realized by
a Bene\v{s} network with $2d-1$ stages \cite{benes}. Each stage flips
one fixed bit as a function of all other bits. For the construction,
pair inputs differing in one chosen bit, and pair outputs the same way.
The permutation gives a degree-two bipartite multigraph on the input
and output pairs. Alternating colors on each cycle, including a doubled
edge, assign one edge of each color to every pair. An input switch and
an output switch along the chosen bit place the colors in separate
classes. Recurse within each class on the remaining $d-1$ bits, with
one switch when $d=1$. The successive target positions are fixed, and
the switch decisions are computable from the permutation table in time
polynomial in $2^d$. Inverting the tables gives the inverse operation.

Align these schedules within each batch, padding shorter schedules by
identities. At one stage keep all active block bits except the current
targets among the controls. Together with the external controls they
use at most $(1-\delta/2)k$ positions. There is at most one target per
block, so a single batch of guarded flips implements the stage.
Multiplying the cost of that batch by $O(1+\log k)$ stages and
$O(1+\log k)$ block batches gives the bound.
\end{proof}
